\chapter{Motivic Invariants of Birational Maps}
\label{chap:overview}

% archon:covers MR4681148MotivicInvariantsOfBirationalMaps.lean MR4681148MotivicInvariantsOfBirationalMaps/Basic.lean

% SOURCE: references/MR4681148-birational-motivic-invariants.tex
% Sections 2.1--2.5, Lemma 2.2, Theorem 2.8, Lemma 2.7, Lemma 2.9, Proposition 2.10.

This chapter formalizes the core construction of Section~2 of Lin--Shinder
\textit{(Motivic invariants of birational maps, arXiv:2207.07389)}: the motivic
invariants \(c\) and \(\tilde{c}\) of birational maps between varieties over an
arbitrary field \(\bk\), together with the key structural results about truncated
Grothendieck groups.

\section{Notation and conventions}

% SOURCE: references/MR4681148-birational-motivic-invariants.tex §1.4

Throughout, \(\bk\) denotes a field.  A \emph{variety over \(\bk\)} is an
irreducible, reduced, separated scheme of finite type over \(\bk\).

% SOURCE: Lin--Shinder arXiv:2207.07389, §1.4 (Notation and conventions)  (read from references/MR4681148-birational-motivic-invariants.tex)
% SOURCE QUOTE: "By a variety over $\k$,
% we mean
% an
% irreducible
% and reduced separated
% scheme
% of finite type over $\k$."
\begin{definition}[Variety]
  \label{def:variety}
  \lean{MR4681148MotivicInvariantsOfBirationalMaps.Basic.TODO.Variety}
  \textit{Source: Lin--Shinder, §1.4.}
  A \emph{variety over \(\bk\)} is an irreducible, reduced, separated scheme of
  finite type over \(\operatorname{Spec}\bk\).  Its \emph{dimension}
  \(\dim X \in \mathbb{N}\) is the Krull dimension.
\end{definition}

% SOURCE: Lin--Shinder arXiv:2207.07389, §1.4 (a)  (read from references/MR4681148-birational-motivic-invariants.tex)
% SOURCE QUOTE: "$\Z[\Bir/\k]$ (resp. $\Z[\Bir_n/\k]$)
% is the free abelian group generated by birational isomorphism classes of varieties
% (resp. varieties of dimension $n$)
% over $\k$.
% By definition
% $\Z[\Bir/\k] = \bigoplus_{n \ge 0}
% \Z[\Bir_n/\k]$. This is the graded Burnside ring of Kontsevich--Tschinkel \cite{KontsevichTschinkel};"
\begin{definition}[Free abelian group on birational classes]
  \label{def:burnside-ring}
  \lean{MR4681148MotivicInvariantsOfBirationalMaps.Basic.TODO.FreeBiratClasses}
  \uses{def:variety}
  \textit{Source: Lin--Shinder, §1.4 (a).}
  Let \(\Bir_n/\bk\) denote the set of birational isomorphism classes of
  \(n\)-dimensional varieties over \(\bk\).  The \emph{free abelian group on
  birational classes in dimension \(n\)} is
  \[
    \ZBir{n} \;:=\; \bigoplus_{[X] \in \Bir_n/\bk} \Z \cdot [X].
  \]
  Its elements are finite formal integer linear combinations of birational
  equivalence classes of \(n\)-dimensional \(\bk\)-varieties.  Setting
  \(\Z[\Bir/\bk] := \bigoplus_{n \geq 0} \ZBir{n}\) gives the graded Burnside ring
  of Kontsevich--Tschinkel.
\end{definition}

\section{The motivic invariants}

% SOURCE: references/MR4681148-birational-motivic-invariants.tex §2.1

% SOURCE: Lin--Shinder arXiv:2207.07389, §2.1  (read from references/MR4681148-birational-motivic-invariants.tex)
% SOURCE QUOTE: "A birational map $\phi : X \dto Y$
% is a morphism $U \to Y$ over $\k$
% defined on a Zariski dense open subscheme
% $U \subset X$
% which is isomorphic onto its image.
% Two birational maps $\phi,\psi : X \dto Y$
% are identical
% if there exists an open
% $U \subset X$ such that $\phi_{|U} = \psi_{|U}$
% as morphisms.
% We can compose birational maps and each
% birational map $\phi$ has an inverse $\phi^{-1}$,
% thereby defining a
% groupoid $\ul{\Bir}/\k$ of birational classes
% consisting of $\k$-varieties as objects
% and birational maps as morphisms."
\begin{definition}[Birational map]
  \label{def:birational-map}
  \lean{MR4681148MotivicInvariantsOfBirationalMaps.Basic.TODO.BirationalMap}
  \uses{def:variety}
  \textit{Source: Lin--Shinder, §2.1.}
  Let \(X\) and \(Y\) be varieties over \(\bk\).  A \emph{birational map}
  \(\phi \colon X \dto Y\) is a morphism \(\phi|_U \colon U \to Y\) defined on a
  Zariski-dense open subscheme \(U \subset X\) that is an open immersion onto its
  image.  Two birational maps are identified when they agree on some common dense
  open.  Birational maps compose (with the natural domain restriction) and every
  birational map \(\phi\) has an inverse \(\phi^{-1}\).  The resulting structure is
  a \emph{groupoid} \(\underline{\Bir}/\bk\) with varieties as objects and
  birational maps as morphisms.
\end{definition}

% SOURCE: Lin--Shinder arXiv:2207.07389, §2.1  (read from references/MR4681148-birational-motivic-invariants.tex)
% SOURCE QUOTE: "We say that $\phi$ is an isomorphism at a point
% $x \in X$ if there exists a Zariski open $U \subset X$
% containing $x$ such that
% $\phi_{|U}$ is an open embedding into $Y$.
% Here, the point $x \in X$ is not required to be closed.
% We define the exceptional set as
% $$\Ex(\phi) \cnec \Set{x \in X | \phi \text{ is not an isomorphism at } x} \subset X.$$
% According to the definition, $\Ex(\phi)$ is a closed
% subset
% containing the indeterminacy locus of $\phi$."
\begin{definition}[Exceptional set]
  \label{def:exceptional-set}
  \lean{MR4681148MotivicInvariantsOfBirationalMaps.Basic.TODO.exceptionalSet}
  \uses{def:birational-map}
  \textit{Source: Lin--Shinder, §2.1.}
  Let \(\phi \colon X \dto Y\) be a birational map.  We say \(\phi\) is an
  \emph{isomorphism at \(x \in X\)} if there exists a Zariski-open \(U \subset X\)
  containing \(x\) such that \(\phi|_U\) is an open immersion into \(Y\).
  The \emph{exceptional set} is
  \[
    \Ex(\phi) \;:=\; \bigl\{\, x \in X \;\big|\; \phi \text{ is not an isomorphism at } x \,\bigr\} \;\subset\; X.
  \]
  By definition, \(\Ex(\phi)\) is a closed subset containing the indeterminacy
  locus of \(\phi\).

  The birational map \(\phi\) restricts to an isomorphism
  \(X \setminus \Ex(\phi) \xrightarrow{\;\sim\;} Y \setminus \Ex(\phi^{-1})\).
\end{definition}

% SOURCE: Lin--Shinder arXiv:2207.07389, §2.1, unnumbered Def (c(φ))  (read from references/MR4681148-birational-motivic-invariants.tex)
% SOURCE QUOTE: "\begin{Def}
% For every birational map $\phi: X \dto Y$ of $\k$-varieties,
% we set
% \begin{equation}\label{eq:def-c}
% c(\phi) \ \ \cnec
% \sum_{\substack{y \in \Ex(\phi^{-1}) \\ \codim_Y y = 1}} [\k(y)]
%  \ \ -  \sum_{\substack{x \in \Ex(\phi) \\ \codim_X x = 1}} [\k(x)] \ \ \in \ \ \Z[\Bir_{n-1}/\k]
%  \end{equation}
% where $\Z[\Bir_{n-1}/\k]$
% is the free abelian group generated by $\Bir_{n-1}/\k$,
% and where the function fields $\k(y)$, $\k(x)$ are identified
% with the corresponding birational classes.
% \end{Def}"
\begin{definition}[Motivic invariant \(c\)]
  \label{def:c-invariant}
  \lean{MR4681148MotivicInvariantsOfBirationalMaps.Basic.TODO.cInvariant}
  \uses{def:exceptional-set, def:burnside-ring}
  \textit{Source: Lin--Shinder, §2.1, unnumbered Def.}
  For a birational map \(\phi \colon X \dto Y\) between \(n\)-dimensional
  \(\bk\)-varieties, define
  \[
    c(\phi) \;:=\;
    \sum_{\substack{y \in \Ex(\phi^{-1}) \\ \codim_Y y = 1}} [\bk(y)]
    \;-\;
    \sum_{\substack{x \in \Ex(\phi) \\ \codim_X x = 1}} [\bk(x)]
    \;\in\; \ZBir{n-1},
  \]
  where the function field \(\bk(y)\) (resp.\ \(\bk(x)\)) of a codimension-one
  point identifies with the corresponding \((n{-}1)\)-dimensional birational class.
  Both sums are finite since \(\Ex(\phi)\) and \(\Ex(\phi^{-1})\) are closed.
\end{definition}

% SOURCE: Lin--Shinder arXiv:2207.07389, §2.1, lem-homc  (read from references/MR4681148-birational-motivic-invariants.tex)
% SOURCE QUOTE: "\begin{lem}\label{lem-homc}
% Given birational maps $\phi : X \dto Y$ and $\psi : Y \dto Z$,
% we have
% $$c(\psi \circ \phi) = c(\phi) + c(\psi).$$
% \end{lem}"
\begin{lemma}[Homomorphism property of \(c\)]
  \label{lem:homc}
  \lean{MR4681148MotivicInvariantsOfBirationalMaps.Basic.TODO.cInvariant_comp}
  \uses{def:c-invariant, thm:def-c, def:pi-n}
  \textit{Source: Lin--Shinder, Lemma 2.2.}

  For birational maps \(\phi \colon X \dto Y\) and \(\psi \colon Y \dto Z\),
  \[
    c(\psi \circ \phi) \;=\; c(\phi) + c(\psi).
  \]
\end{lemma}
% SOURCE QUOTE PROOF: "Lemma~\ref{lem-homc} has a natural proof
% based on motivic cut-and-paste,
% and we postpone the proof until Subsection
% \ref{ssec-leftK},
% see Theorem \ref{thm-Defc}."
\begin{proof}
  This follows from the existence and uniqueness of \(\tilde{c}\) in
  \cref{thm:def-c}: since \(\tilde{c}\) satisfies condition (ii) by construction,
  composing with the map \(\pi_{n-1} \colon \KVar{n-1} \to \ZBir{n-1}\) (which is
  also additive) yields the additivity of \(c = \pi_{n-1} \circ \tilde{c}\).
\end{proof}

\subsection{Truncated Grothendieck groups \texorpdfstring{\(\KVar{n}\)}{K0(Var<=n/k)}}

% SOURCE: references/MR4681148-birational-motivic-invariants.tex §2.3

% SOURCE: Lin--Shinder arXiv:2207.07389, §2.3  (read from references/MR4681148-birational-motivic-invariants.tex)
% SOURCE QUOTE: "We write $\Kt_0(\Var^{\le n}/\k)$ for the Grothendieck
% group of varieties of dimension $\le n$: the generators
% of $\Kt_0(\Var^{\le n}/\k)$
% are
% classes $[X]_n$
% of schemes of finite type $X/\k$ with $\dim(X) \le n$
% and the relations
% are generated by
% \begin{equation}\label{eq:cut-and-paste}
% [X]_n = [U]_n + [Z]_n
% \end{equation}
% for every open $U \subset X$ with closed complement $Z$.
% We have a sequence of natural maps
% \begin{equation}\label{eqn-iota}
%     \cdots \to \Kt_0(\Var^{\le n-1}/\k) \xto{\iota_{n-1}} \Kt_0(\Var^{\le n}/\k) \xto{\iota_{n}} \Kt_0(\Var^{\le n+1}/\k) \to \cdots.
% \end{equation}"
\begin{definition}[Truncated Grothendieck group]
  \label{def:truncated-grothendieck}
  \lean{MR4681148MotivicInvariantsOfBirationalMaps.Basic.TODO.TruncatedGrothendieckGroup}
  \uses{def:variety}
  \textit{Source: Lin--Shinder, §2.3.}
  Let \(n \geq 0\).  The \emph{truncated Grothendieck group} \(\KVar{n}\) is the
  abelian group with:
  \begin{itemize}
    \item \emph{Generators}: isomorphism classes \([X]_n\) of \(\bk\)-schemes of
      finite type with \(\dim X \leq n\).
    \item \emph{Relations}: \([X]_n = [U]_n + [Z]_n\) for every open subscheme
      \(U \subset X\) with closed complement \(Z = X \setminus U\).
  \end{itemize}
  The natural map \(\iota_{n-1} \colon \KVar{n-1} \to \KVar{n}\) (sending
  \([X]_{n-1} \mapsto [X]_n\)) is the inclusion induced by the identity on
  generators.  The colimit \(\varinjlim_n \KVar{n}\) is the underlying abelian
  group of the Grothendieck ring of varieties \(\mathrm{K}_0(\Var/\bk)\).
\end{definition}

% SOURCE: Lin--Shinder arXiv:2207.07389, §2.3, eq.~(def-pi)  (read from references/MR4681148-birational-motivic-invariants.tex)
% SOURCE QUOTE: "where
% we define
% $\pi_n$
% on
% the classes of
% finite type schemes
% of dimension $\le n$ by
% \begin{equation}
% \label{eq:def-pi}
% \pi_n([X]_n) = \left\{\begin{array}{cc}
%      \; [X_1] + \cdots + [X_r], & \dim(X) = n \\
%      \; 0, & \dim(X) < n \\
% \end{array}\right.,
% \end{equation}
% where $X_1,\ldots,X_r$ are the irreducible components of $X$ of dimension $n$."
\begin{definition}[Projection \(\pi_n\)]
  \label{def:pi-n}
  \lean{MR4681148MotivicInvariantsOfBirationalMaps.Basic.TODO.piMap}
  \uses{def:truncated-grothendieck, def:burnside-ring}
  \textit{Source: Lin--Shinder, §2.3, eq.~(2.5).}
  Define the group homomorphism
  \[
    \pi_n \colon \KVar{n} \;\to\; \ZBir{n}
  \]
  on generators by
  \[
    \pi_n([X]_n) \;:=\;
    \begin{cases}
      [X_1] + \cdots + [X_r], & \dim X = n \\
      0, & \dim X < n
    \end{cases}
  \]
  where \(X_1, \ldots, X_r\) are the \(n\)-dimensional irreducible components of
  \(X\).  This is well-defined: if \(U \subset X\) is open with closed complement
  \(Z\), the \(n\)-dimensional components of \(X\) are distributed between \(U\) and
  \(Z\) (those meeting both \(U\) and \(Z\) lie entirely in \(U\) and their boundary
  in \(Z\) has dimension \(< n\)), so \(\pi_n([X]_n) = \pi_n([U]_n) + \pi_n([Z]_n)\).
\end{definition}

% SOURCE: Lin--Shinder arXiv:2207.07389, §2.3, eq.~(K0-seq)  (read from references/MR4681148-birational-motivic-invariants.tex)
% SOURCE QUOTE: "For the cokernel, we have the exact sequence
% \begin{equation}\label{eq:K0-seq}
% \Kt_0(\Var^{\le n-1}/\k) \overset{\iota_{n-1}}\to \Kt_0(\Var^{\le n}/\k) \overset{\pi_n}\to \Z[\Bir_n/\k] \to 0,
% \end{equation}"
\begin{lemma}[Exact sequence]
  \label{lem:exact-seq}
  \lean{MR4681148MotivicInvariantsOfBirationalMaps.Basic.TODO.exactSeq}
  \uses{def:truncated-grothendieck, def:pi-n, lem:kernelXU}
  \textit{Source: Lin--Shinder, §2.3, eq.\ (2.3).}

  For every \(n\), the sequence
  \[
    \KVar{n-1} \xrightarrow{\;\iota_{n-1}\;} \KVar{n} \xrightarrow{\;\pi_n\;} \ZBir{n} \;\to\; 0
  \]
  is exact.
\end{lemma}
% SOURCE QUOTE PROOF: "For the cokernel, we have the exact sequence
% \begin{equation}\label{eq:K0-seq}
% \Kt_0(\Var^{\le n-1}/\k) \overset{\iota_{n-1}}\to \Kt_0(\Var^{\le n}/\k) \overset{\pi_n}\to \Z[\Bir_n/\k] \to 0,
% \end{equation}
% where
% we define
% $\pi_n$
% on
% the classes of
% finite type schemes
% of dimension $\le n$ by
% \begin{equation}
% \label{eq:def-pi}
% \pi_n([X]_n) = \left\{\begin{array}{cc}
%      \; [X_1] + \cdots + [X_r], & \dim(X) = n \\
%      \; 0, & \dim(X) < n \\
% \end{array}\right.,
% \end{equation}
% where $X_1,\ldots,X_r$ are the irreducible components of $X$ of dimension $n$.
% [...]
% See Lemma \ref{lem:kernelXY}
% for a reformulation of the statement using the motivic invariant
% $\ti{c}$."
\begin{proof}
  \(\pi_n \circ \iota_{n-1} = 0\) since all generators of \(\KVar{n-1}\) have
  dimension \(< n\) and are sent to \(0\) by \(\pi_n\).  Surjectivity of \(\pi_n\):
  every birational class \([X] \in \Bir_n/\bk\) is realized by an irreducible
  \(n\)-dimensional variety, and \(\pi_n([X]_n) = [X]\).  Exactness at \(\KVar{n}\)
  follows from the presentation of \cref{lem:kernelXU}: every element of
  \(\ker(\pi_n)\) is a combination of classes of dimension \(< n\) after
  cut-and-paste, hence lies in the image of \(\iota_{n-1}\).
\end{proof}

% SOURCE: Lin--Shinder arXiv:2207.07389, §2.3, lem:kernelXU  (read from references/MR4681148-birational-motivic-invariants.tex)
% SOURCE QUOTE: "\begin{lemma}\label{lem:kernelXU}
% Let $\Var^n/\k$ denote the set of isomorphism
% classes of $n$-dimensional $\k$-varieties.
% Then the natural homomorphism
% \begin{equation}\label{eq:plusprojectionU}
% \Kt_0(\Var^{\le n-1}/\k) \oplus \Z[\Var^n/\k] \to \Kt_0(\Var^{\le n}/\k)
% \end{equation}
% is surjective and its kernel
% admits the following presentation:
% \begin{equation}\label{eq:kernel-XU}
% \langle \ \left( [X \setminus U]_{n-1},
% - [X] + [U] \right) \ \rangle_{U \subset X},
% \end{equation}
% for all $n$-dimensional
% varieties $X$
% and open subsets $U \subset X$.
% \end{lemma}"
\begin{lemma}[Presentation of the kernel]
  \label{lem:kernelXU}
  \lean{MR4681148MotivicInvariantsOfBirationalMaps.Basic.TODO.lem_kernelXU}
  \uses{def:truncated-grothendieck}
  \textit{Source: Lin--Shinder, Lemma 2.7.}

  Let \(\mathrm{Var}^n/\bk\) be the set of isomorphism classes of
  \(n\)-dimensional \(\bk\)-varieties.  The natural group homomorphism
  \[
    \beta \colon \KVar{n-1} \oplus \Z[\mathrm{Var}^n/\bk] \;\to\; \KVar{n},
    \quad
    ([Z]_{n-1}, [X]) \;\mapsto\; [Z]_n + [X]_n,
  \]
  is surjective.  Its kernel is the subgroup generated by all elements of the form
  \[
    \bigl([X \setminus U]_{n-1},\; -[X] + [U]\bigr)
  \]
  as \(U\) ranges over all open subsets of all irreducible \(n\)-dimensional
  varieties \(X\) over \(\bk\).  Equivalently, \(\beta\) induces an isomorphism
  \[
    \frac{\KVar{n-1} \oplus \Z[\mathrm{Var}^n/\bk]}
         {\bigl\langle ([X \setminus U]_{n-1}, -[X]+[U]) \bigr\rangle_{U \subset X}}
    \;\xrightarrow{\;\sim\;}\; \KVar{n}.
  \]
\end{lemma}
% SOURCE QUOTE PROOF: "It is clear from definitions that
% \eqref{eq:plusprojectionU}
% is surjective and that
% its kernel contains
% \eqref{eq:kernel-XU},
% so that we have a well-defined surjective homomorphism
% \begin{equation}\label{eq:plusprojectionU-new}
% \frac{\Kt_0(\Var^{\le n-1}/\k) \oplus \Z[\Var^n/\k]}{\langle \ \left( [X \setminus U]_{n-1},
% - [X] + [U] \right) \ \rangle_{U \subset X}}
% \overset{\beta}\longrightarrow
% \Kt_0(\Var^{\le n}/\k).
% \end{equation}
% To show that \eqref{eq:plusprojectionU-new} is also
% injective,
% we construct the inverse homomorphism $\alpha$ of $\gb$.
% First we define $\alpha(X)$ for a scheme $X/\k$
% of finite type of dimension $\le n$.
% Since the class $[X]_n$ only depends on the reduced structure,
% we can assume that $X$ is reduced.
% We stratify $X$ into
% irreducible varieties which are locally closed in $X$ and
% let $X_1, \dots, X_r$
% be its $n$-dimensional strata
% ($r = 0$ if $\dim(X) < n$).
% The element
% \[
% \alpha(X) :=
% \left([X \setminus \bigsqcup_{i=1}^r X_i ]_{n-1}, \sum_{i=1}^r [X_i]\right)
% \]
% is well-defined
% (that is, independent of the choice of stratification of $X$)
% when
% considered in the left-hand side of \eqref{eq:plusprojectionU-new}.
% We verify that
% $\alpha(X) = \alpha(V) + \alpha(X\setminus V)$
% for every open subscheme $V \subset X$, so
% that $\alpha$ defines a homomorphism
% from $\Kt_0(\Var^{\le n}/\k)$.
% By construction $\alpha$ and $\beta$ are mutually inverse."
\begin{proof}
  Surjectivity is clear: any \([X]_n \in \KVar{n}\) is the image of \((0, [X])\) if
  \(\dim X = n\), or \(([X]_{n-1}, 0)\) if \(\dim X < n\).  The given elements are
  in the kernel because \([X]_n = [U]_n + [X \setminus U]_n\) by the cut-and-paste
  relation.

  For injectivity of the induced map, we construct the inverse \(\alpha\) of
  \(\beta\) as follows.  Given a scheme \(X\) of dimension \(\leq n\), stratify \(X\)
  into locally closed irreducible strata; let \(X_1, \ldots, X_r\) be the
  \(n\)-dimensional strata.  Set
  \[
    \alpha([X]_n) \;:=\;
    \Bigl([X \setminus \bigsqcup_i X_i]_{n-1},\; \sum_i [X_i]\Bigr).
  \]
  This is independent of the choice of stratification (when taken in the quotient
  by the given relations), and satisfies \(\alpha([U]_n + [Z]_n) = \alpha([X]_n)\)
  for open \(U \subset X\) with closed complement \(Z\).  One checks that
  \(\alpha \circ \beta = \mathrm{id}\) and \(\beta \circ \alpha = \mathrm{id}\).
\end{proof}

\subsection{Invariant \texorpdfstring{\(\tilde{c}(\phi)\)}{c-tilde(phi)} with values in \texorpdfstring{\(\KVar{n-1}\)}{K0(Var<=n-1/k)}}

% SOURCE: references/MR4681148-birational-motivic-invariants.tex §2.4

% SOURCE: Lin--Shinder arXiv:2207.07389, §2.4, eq.~(def-ct)  (read from references/MR4681148-birational-motivic-invariants.tex)
% SOURCE QUOTE: "\begin{equation}\label{eq:def-ct}
%     \ti{c}(\phi) = \ti{c}(j) - \ti{c}(i) = [Y \setminus j(U)]_{n-1} - [X \setminus U]_{n-1}
% \end{equation}"
\begin{definition}[Enhanced invariant \(\tilde{c}\)]
  \label{def:c-tilde}
  \lean{MR4681148MotivicInvariantsOfBirationalMaps.Basic.TODO.cTildeInvariant}
  \uses{def:birational-map, def:truncated-grothendieck, thm:def-c}
  \textit{Source: Lin--Shinder, §2.4, eq.~(2.7).}
  For a birational map \(\phi \colon X \dto Y\) between \(n\)-dimensional
  \(\bk\)-varieties, choose any open \(U \subset X\) on which \(\phi|_U \colon U \to Y\)
  is an open immersion.  Define
  \[
    \tilde{c}(\phi) \;:=\; [Y \setminus \phi(U)]_{n-1} - [X \setminus U]_{n-1}
    \;\in\; \KVar{n-1}.
  \]
  By \cref{thm:def-c} below, this is independent of the choice of \(U\) and
  satisfies the homomorphism property.
\end{definition}

% SOURCE: Lin--Shinder arXiv:2207.07389, §2.4, thm-Defc  (read from references/MR4681148-birational-motivic-invariants.tex)
% SOURCE QUOTE: "\begin{theorem}\label{thm-Defc}
% There exists a unique assignment
% \[
% \ti{c}: \Bir(X,Y) \to \Kt_0(\Var^{\le{n-1}}/\k)
% \]
% defined on the birational maps
% between $n$-dimensional varieties over $\k$ which
% satisfies the following two properties.
% \begin{enumerate}
%     \item If $i: U \hto X$ is the inclusion of an open subset
%     then $\ti{c}(i) = [X \setminus U]_{n-1}$.
%     \item For every $\phi \in \Bir(X, Y)$, $\psi \in \Bir(Y,Z)$,
%     we have $\ti{c}(\psi \circ \phi) = \ti{c}(\phi) + \ti{c}(\psi)$.
% \end{enumerate}
% Furthermore, we have $c(\phi) = \pi_{n-1}(\ti{c}(\phi))$,
% where $c(\phi)$ is defined by \eqref{eq:def-c} and
% $\pi_{n-1}$ is defined by \eqref{eq:def-pi}.
% \end{theorem}"
\begin{theorem}[Existence and uniqueness of \(\tilde{c}\)]
  \label{thm:def-c}
  \lean{MR4681148MotivicInvariantsOfBirationalMaps.Basic.TODO.thm_DefC}
  \uses{def:birational-map, def:truncated-grothendieck, def:c-invariant, def:pi-n}
  \textit{Source: Lin--Shinder, Theorem 2.8.}

  There exists a unique assignment
  \[
    \tilde{c} \colon \Bir(X, Y) \;\to\; \KVar{n-1}
  \]
  defined on birational maps between \(n\)-dimensional \(\bk\)-varieties satisfying:
  \begin{enumerate}
    \item[(i)] If \(i \colon U \hto X\) is the inclusion of an open subset, then
      \(\tilde{c}(i) = [X \setminus U]_{n-1}\).
    \item[(ii)] For \(\phi \in \Bir(X,Y)\) and \(\psi \in \Bir(Y,Z)\),
      \(\tilde{c}(\psi \circ \phi) = \tilde{c}(\phi) + \tilde{c}(\psi)\).
  \end{enumerate}
  Moreover, \(c(\phi) = \pi_{n-1}(\tilde{c}(\phi))\), where \(c\) is the invariant
  of \cref{def:c-invariant}.
\end{theorem}
% SOURCE QUOTE PROOF: "Let $\phi \in \Bir(X, Y)$, then there is an open subset $i: U \subset X$
% such that $j = \phi|_U$ is an open embedding. We have $\phi = j \circ i^{-1}$
% so that properties (1) and (2) determine $\ti{c}(\phi)$ to be
% \begin{equation}\label{eq:def-ct}
%     \ti{c}(\phi) = \ti{c}(j) - \ti{c}(i) = [Y \setminus j(U)]_{n-1} - [X \setminus U]_{n-1}
% \end{equation}
% This shows that $\ti{c}$ is unique.
% To prove the existence,
% we first check that \eqref{eq:def-ct} is well-defined, that is $\ti{c}(\phi)$
% is independent of
% the choice of $U \subset X$.
% [...]
% To check (2), let $\phi \in \Bir(X, Y)$ and $\psi \in \Bir(Y,Z)$,
% and choose $U \subset X$ such that both $j = \phi|_U$ and
% $j'= \psi|_{j(U)}$ are open embeddings. Then
% $$\ti{c}(\psi) + \ti{c}(\phi)
% = [Z \setminus j'(j(U))]_{n-1} - [Y \setminus j(U)]_{n-1}  +
% [Y \setminus j(U)]_{n-1} - [X \setminus U]_{n-1} = \ti{c}(\psi
% \circ \phi).$$
% For the final claim, by \eqref{eq:iso-ex} we can factorize $\phi$ through
% the inclusions
% \[
% X \hlto X \setminus \Ex(\phi) \simeq Y \setminus \Ex(\phi^{-1}) \hto Y,
% \]
% so that by (1) and (2) we have
% $\ti{c}(\phi) = [\Ex(\phi^{-1})]_{n-1} - [\Ex(\phi)]_{n-1}$.
% It follows from the definitions that $\pi_{n-1}(\ti{c}(\phi)) = c(\phi)$."
\begin{proof}
  \textbf{Uniqueness.}  Every birational map \(\phi \colon X \dto Y\) factors as
  \(\phi = j \circ i^{-1}\) where \(i \colon U \hto X\) and \(j = \phi|_U \colon U \hto Y\)
  are open immersions.  Conditions (i) and (ii) then force
  \[
    \tilde{c}(\phi) = \tilde{c}(j) - \tilde{c}(i)
    = [Y \setminus j(U)]_{n-1} - [X \setminus U]_{n-1}.
  \]

  \textbf{Well-definedness.}  If \(U' \subset U\) is another open domain on which
  \(\phi|_{U'}\) is an open immersion, then using the cut-and-paste relations in
  \(\KVar{n-1}\):
  \begin{align*}
    [Y \setminus \phi(U')]_{n-1} - [X \setminus U']_{n-1}
    &= \bigl([Y \setminus \phi(U)]_{n-1} + [\phi(U) \setminus \phi(U')]_{n-1}\bigr)
      - \bigl([X \setminus U]_{n-1} + [U \setminus U']_{n-1}\bigr) \\
    &= [Y \setminus \phi(U)]_{n-1} - [X \setminus U]_{n-1},
  \end{align*}
  since \(\phi|_U\) is an isomorphism \(U \setminus U' \xrightarrow{\sim} \phi(U) \setminus \phi(U')\).

  \textbf{Homomorphism property.}  For \(\phi \in \Bir(X,Y)\), \(\psi \in \Bir(Y,Z)\),
  choose \(U \subset X\) such that both \(j = \phi|_U\) and \(\psi|_{j(U)}\) are
  open immersions.  Then
  \begin{align*}
    \tilde{c}(\psi) + \tilde{c}(\phi)
    &= [Z \setminus \psi(j(U))]_{n-1} - [Y \setminus j(U)]_{n-1}
      + [Y \setminus j(U)]_{n-1} - [X \setminus U]_{n-1} \\
    &= [Z \setminus (\psi \circ \phi)(U)]_{n-1} - [X \setminus U]_{n-1}
    = \tilde{c}(\psi \circ \phi).
  \end{align*}

  \textbf{Relation \(c = \pi_{n-1} \circ \tilde{c}\).}  Factor \(\phi\) through
  \(X \setminus \Ex(\phi) \xrightarrow{\sim} Y \setminus \Ex(\phi^{-1})\).  By
  (i) and (ii),
  \[
    \tilde{c}(\phi) = [\Ex(\phi^{-1})]_{n-1} - [\Ex(\phi)]_{n-1}.
  \]
  Applying \(\pi_{n-1}\) picks out the codimension-one components and reproduces
  the definition of \(c(\phi)\).
\end{proof}

\subsection{Invariant \texorpdfstring{\(\tilde{c}\)}{c-tilde} and the structure of the Grothendieck rings}

% SOURCE: references/MR4681148-birational-motivic-invariants.tex §2.5

% SOURCE: Lin--Shinder arXiv:2207.07389, §2.5, lem:kernelXY  (read from references/MR4681148-birational-motivic-invariants.tex)
% SOURCE QUOTE: "\begin{lemma}\label{lem:kernelXY}
% The kernel of \eqref{eq:plusprojectionU}
% admits the following presentation:
% \begin{equation}\label{eq:kernel-XY}
% \langle \ \left( \ti{c}(\phi),
% -[Y] + [X] \right) \ \rangle_{\phi \in \Bir(X,Y)}
% \end{equation}
% for all birational isomorphisms $\phi$
% between all irreducible $n$-dimensional
% varieties $X$, $Y$.
% \end{lemma}"
\begin{lemma}[Kernel presentation via \(\tilde{c}\)]
  \label{lem:kernelXY}
  \lean{MR4681148MotivicInvariantsOfBirationalMaps.Basic.TODO.lem_kernelXY}
  \uses{lem:kernelXU, thm:def-c, def:c-tilde}
  \textit{Source: Lin--Shinder, Lemma 2.9.}

  The kernel of \(\beta \colon \KVar{n-1} \oplus \Z[\mathrm{Var}^n/\bk] \to \KVar{n}\)
  (from \cref{lem:kernelXU}) admits the alternative presentation
  \[
    \ker(\beta)
    = \Bigl\langle \bigl(\tilde{c}(\phi),\; -[Y] + [X]\bigr) \Bigr\rangle_{\phi \in \Bir(X,Y)}
  \]
  as \(\phi\) ranges over all birational isomorphisms between all irreducible
  \(n\)-dimensional varieties \(X\) and \(Y\).
\end{lemma}
% SOURCE QUOTE PROOF: "The subgroup \eqref{eq:kernel-XU} can be rewritten as
% \begin{equation}\label{eq:kernel-UX}
% \langle \ \left( \ti{c}(j:U \hto X),
% - [X] + [U] \right) \ \rangle_{U \subset X}
% \end{equation}
% where
% $U \subset X$ runs over
% open subsets of all irreducible $n$-dimensional
% varieties $X$.
% It remains to show that the subgroups \eqref{eq:kernel-UX}
% and \eqref{eq:kernel-XY} coincide.
% It is clear that \eqref{eq:kernel-UX} is a subgroup of
% \eqref{eq:kernel-XY}. Conversely, every element of
% \eqref{eq:kernel-XY} can be written as a combination
% of elements
% from \eqref{eq:kernel-UX} after decomposing
% $\phi$ as a composition of open embeddings and their inverses."
\begin{proof}
  From \cref{lem:kernelXU}, the kernel is generated by \(([X \setminus U]_{n-1}, -[X]+[U])\)
  for open \(U \subset X\).  Identifying \(\tilde{c}(U \hto X) = [X \setminus U]_{n-1}\)
  (condition (i) of \cref{thm:def-c}), each such generator is
  \((\tilde{c}(i), -[X]+[U])\) where \(i \colon U \hto X\).  This shows the
  generator set from \cref{lem:kernelXU} is contained in the one for the current lemma.
  Conversely, every birational map \(\phi \colon X \dto Y\) decomposes as a
  composition of open immersions and their inverses; the corresponding
  \((\tilde{c}(\phi), -[Y]+[X])\) is a combination of generators of the first type
  by the homomorphism property of \(\tilde{c}\).
\end{proof}

% SOURCE: Lin--Shinder arXiv:2207.07389, §2.5, prop:ker-iota  (read from references/MR4681148-birational-motivic-invariants.tex)
% SOURCE QUOTE: "\begin{proposition}\label{prop:ker-iota}
% We have
% \begin{equation}\label{eq:ker-iota}
% \Ker(\iota_{n-1}) = \sum_{X \in \Bir_n/\k} \ti{c}(\Bir(X))
% \end{equation}
% so that there is an exact sequence
% \begin{equation}\label{eq:seq-K0}
% 0 \to \Image(\ti{c})_n \to \Kt_0(\Var^{\le n-1}/\k) \overset{\iota_{n-1}}\to \Kt_0(\Var^{\le n}/\k) \overset{\pi_n}\to \Z[\Bir_n/\k] \to 0,
% \end{equation}
% where we write $\Image(\ti{c})_n$ for the right hand side
% of \eqref{eq:ker-iota}.
% \end{proposition}"
\begin{proposition}[Kernel of \(\iota_{n-1}\)]
  \label{prop:ker-iota}
  \lean{MR4681148MotivicInvariantsOfBirationalMaps.Basic.TODO.prop_ker_iota}
  \uses{lem:kernelXY, thm:def-c, def:truncated-grothendieck}
  \textit{Source: Lin--Shinder, Proposition 2.10.}

  We have
  \[
    \ker(\iota_{n-1}) \;=\; \sum_{X \in \Bir_n/\bk} \tilde{c}(\Bir(X)),
  \]
  and there is an exact sequence
  \[
    0 \;\to\; \Image(\tilde{c})_n \;\to\; \KVar{n-1}
      \;\xrightarrow{\;\iota_{n-1}\;}\; \KVar{n}
      \;\xrightarrow{\;\pi_n\;}\; \ZBir{n} \;\to\; 0,
  \]
  where \(\Image(\tilde{c})_n := \sum_{X \in \Bir_n/\bk} \tilde{c}(\Bir(X))\).
\end{proposition}
% SOURCE QUOTE PROOF: "$\Ker(\iota_{n-1})$ coincides with kernel of $\eqref{eq:plusprojectionU}$
% intersected with $\Kt_0(\Var^{n-1}/\k)$.
% Hence by Lemma \ref{lem:kernelXY} every element in
% $\Ker(\iota_{n-1})$ has the form
% \begin{equation}\label{eq:sum-tcphi}
% \sum_{i=1}^r \tc(\phi_i)
% \end{equation}
% with $\phi_i: X_i \dto Y_i$
% and $\sum_{i=1}^r ([X_i] - [Y_i]) = 0 \in \Z[\Var^n/\k]$.
% Since the latter is a free abelian group on the isomorphism
% classes of $n$-dimensional irreducible varieties,
% there is a permutation $\sigma \in S_r$ such that
% $Y_i \simeq X_{\sigma(i)}$.
% Let $i \in \{1, \dots, r\}$ and let $\ell$ be the length
% of $\sigma$-orbit of $i$. Then we have a composition
% \[
% \psi_i: X_i \overset{\phi_i}{\dto} X_{\sigma(i)} \overset{\phi_{\sigma(i)}}{\dto} X_{\sigma^2(i)} \overset{\phi_{\sigma^2(i)}}{\dto} \dots
% \overset{\phi_{\sigma^{\ell-1}(i)}}{\dto}
% X_{\sigma^\ell(i)}  =  X_i.
% \]
% This allows to rewrite \eqref{eq:sum-tcphi}
% as a sum of $\tc(\psi_i)$, with $\psi_i \in \Bir(X_i)$,
% and $i$ running over a set of representative classes
% for orbits of $\sigma$ on $\{1, \dots, r\}$."
\begin{proof}
  \uses{lem:exact-seq, def:pi-n}
  The kernel of \(\beta\) from \cref{lem:kernelXU} intersected with \(\KVar{n-1}\)
  (i.e.\ the first factor) is exactly \(\ker(\iota_{n-1})\).  By \cref{lem:kernelXY},
  the intersection is generated by \(\tilde{c}(\phi)\) subject to the constraint
  \([Y] = [X]\) in \(\Z[\mathrm{Var}^n/\bk]\).  Since this free abelian group has
  isomorphism classes as a basis, the constraint \([Y] = [X]\) forces \(Y \cong X\)
  (isomorphic as varieties, not just birationally equivalent).  Composing the
  individual birational automorphisms (possibly from different orbits) into
  self-birational maps of the same variety, one can rewrite any element of
  \(\ker(\iota_{n-1})\) as a sum \(\sum_i \tilde{c}(\psi_i)\) with
  \(\psi_i \in \Bir(X_i)\), proving the stated description of the kernel.

  Exactness at \(\KVar{n-1}\) follows from the kernel description above.  Exactness
  at \(\KVar{n}\) and surjectivity of \(\pi_n\) are from \cref{lem:exact-seq}.
\end{proof}
